\begin{frame}{Problem setup and design requirements}
\protect\phantomsection\label{problem-setup-and-design-requirements}
\begin{block}{Problem Understanding}
\protect\phantomsection\label{problem-understanding}
The system consists of an inverted pendulum mounted to a motorized cart.

\textbf{Features of this system:}

\begin{enumerate}
\item
  It is unstable without control, i.e., the pendulum will simply fall
  over if the cart isn't moved to balance it.
\item
  The dynamics of the system are nonlinear.
\end{enumerate}

\textbf{Objective:}

The objective of the control system is to balance the inverted pendulum
by applying a

force to the cart that the pendulum is attached to.

\textbf{Assumption:}

We consider a two-dimensional problem where the pendulum is constrained
to move in

the vertical plane. For this system:

\begin{itemize}
\item
  The control input is the force \(F\) that moves the cart horizontally.
\item
  The outputs are the angular position of the pendulum \(\theta\) and
  the horizontal position of the cart \(x\).
\end{itemize}

\begin{figure}[H]
\centering
\includegraphics[width=0.5\linewidth,height=\textheight,keepaspectratio]{plots/system.png}
\caption{Inverted Pendulum}\label{fig:histogram}
\end{figure}
\end{block}

\begin{block}{System Parameters}
\protect\phantomsection\label{system-parameters}
For this example, let's assume the following quantities:

\begin{itemize}
\item
  \(M\) Mass of the cart \(0.5 \, \text{kg}\)
\item
  \(m\) Mass of the pendulum \(0.2 \, \text{kg}\)
\item
  \(b\) Coefficient of friction for cart \(0.1 \, \text{N/m/sec}\)
\item
  \(l\) Length to pendulum center of mass \(0.3 \, \text{m}\)
\item
  \(I\) Mass moment of inertia of the pendulum
  \(0.006 \, \text{kg}\cdot\text{m}^2\)
\item
  \(F\) Force applied to the cart
\item
  \(x\) Cart position coordinate
\item
  \(\theta\) Pendulum angle from vertical (down)
\end{itemize}
\end{block}

\begin{block}{Design Requirements}
\protect\phantomsection\label{design-requirements}
We will design a controller to restore the pendulum to a vertically
upward position after it has experienced an impulsive ``bump'' to the
cart. Specifically, the design criteria are that the pendulum return to
its upright position within 5 seconds and that the pendulum never move
more than 0.05 radians away from vertical after being disturbed by an
impulse of magnitude \(1 \, \text{Nsec}\). The pendulum will initially
begin in the vertically upward equilibrium, \(\theta = \pi\).

\begin{block}{SISO Design Requirements}
\protect\phantomsection\label{siso-design-requirements}
\begin{itemize}
\item
  Settling time for \(\theta\) of less than 5 seconds
\item
  Pendulum angle \(\theta\) never more than \(0.05\) radians from the
  vertical
\end{itemize}

In our case, the inverted pendulum system is single-input, multi-output
(SIMO). Therefore, for the state-space section of the inverted pendulum,
we will attempt to control both the pendulum's angle and the cart's
position.

For this example, We will command a 0.2-meter step in the cart's desired
position.
\end{block}

\begin{block}{State-Space Design Requirements}
\protect\phantomsection\label{state-space-design-requirements}
\begin{itemize}
\item
  Settling time for \(x\) and \(\theta\) of less than 5 seconds
\item
  Rise time for \(x\) of less than 0.5 seconds
\item
  Pendulum angle \(\theta\) never more than \(20^{\circ}\) (\(0.35\)
  radians) from the vertical
\item
  Steady-state error of less than 2\% for \(x\) and \(\theta\)
\end{itemize}
\end{block}
\end{block}
\end{frame}

\begin{frame}{Force analysis and system equations}
\protect\phantomsection\label{force-analysis-and-system-equations}
\begin{figure}[H]
\centering
\includegraphics[width=0.4\linewidth,height=\textheight,keepaspectratio]{plots/forces.png}
\caption{Forces on Inverted Pendulum}\label{fig:histogram}
\end{figure}

\begin{block}{Equations of Motion}
\protect\phantomsection\label{equations-of-motion}
\begin{block}{Pendulum}
\protect\phantomsection\label{pendulum}
\begin{figure}[H]
\centering
\includegraphics[width=0.2\linewidth,height=\textheight,keepaspectratio]{plots/pendulum.png}
\caption{Pendulum}\label{fig:histogram}
\end{figure}

The coordinates of the center of mass are given by: \begin{equation}
(x_c, y_c) \equiv \left( x + l \sin \theta, \, -l \cos \theta \right)
\end{equation}

Velocities: \begin{equation}
(v_x, v_y) \equiv \left( \dot{x} + l \cos \theta \, \dot{\theta}, \; l \sin \theta \, \dot{\theta} \right)
\end{equation}

Accelerations: \begin{equation}
(a_x, a_y) = \left( \ddot{x} - l \sin \theta \, \dot{\theta}^2 + l \cos \theta \, \ddot{\theta}, \;
l \cos \theta \, \dot{\theta}^2 + l \sin \theta \, \ddot{\theta} \right)
\end{equation}

Force equations for the pendulum: \begin{align}
N &= m \ddot{x} \\
P - mg &= m \ddot{y}
\end{align}

Substituting accelerations: \begin{align}
N &= m \left( \ddot{x} - l \sin \theta \, \dot{\theta}^2 + l \cos \theta \, \ddot{\theta} \right) \\
P - mg &= m \left( l \cos \theta \, \dot{\theta}^2 + l \sin \theta \, \ddot{\theta} \right)
\end{align}
\end{block}

\begin{block}{Cart}
\protect\phantomsection\label{cart}
\begin{figure}[H]
\centering
\includegraphics[width=0.3\linewidth,height=\textheight,keepaspectratio]{plots/cart.png}
\caption{Cart}\label{fig:histogram}
\end{figure}

Force balance on the cart: \begin{equation}
F - N - b \ddot{x} = M \ddot{x}
\end{equation}

Substituting \(N\): \begin{equation}
M \ddot{x} = F - b \ddot{x} - m \left( \ddot{x} - l \sin \dot{\theta} \, \dot{\theta}^2 + l \cos \theta \, \ddot{\theta} \right)
\end{equation}

Rearranging: \begin{equation}
(M + m) \ddot{x} + b \ddot{x} = F + m \left( l \sin \theta \, \dot{\theta}^2 - l \cos \theta \, \ddot{\theta} \right)
\end{equation}

Thus, the final cart equation is: \begin{equation}
F = (M + m) \ddot{x} + b \dot{x} + m \left( l \cos \theta \, \ddot{\theta} - l \sin \theta \, \dot{\theta}^2 \right)
\label{eq:cart}
\end{equation}
\end{block}

\begin{block}{Rotation (Moment) Equation}
\protect\phantomsection\label{rotation-moment-equation}
\begin{figure}[H]
\centering
\includegraphics[width=0.3\linewidth,height=\textheight,keepaspectratio]{plots/rotation.png}
\caption{Rotation Forces}\label{fig:histogram}
\end{figure}

According to Newton--Euler formulation:

\begin{equation}
\text{Torque} = \text{rate of change of angular momentum about pivot}
\end{equation}

Thus, the torque balance is: \begin{equation}
\sum \tau_{\text{pivot}} = I_p \ddot{\theta}
\end{equation}

where the moment of inertia about the pivot is obtained using the
parallel axis theorem: \begin{equation}
I_p = I + m l^2
\end{equation}

Torque balance about the pivot: \begin{equation}
P \sin \theta + N \cos \theta - mg \sin \theta 
= m l \ddot{\theta} + m \ddot{x} \cos \theta
\end{equation}

Moment about the center of mass (COM): \begin{equation}
-(P \sin \theta) l - (N \cos \theta) l = I \ddot{\theta}
\end{equation}

Combining both and eliminating \(P, N\): \begin{equation}
- m g l \sin \theta = m l \ddot{\theta} + m \ddot{x} l \cos \theta + I \ddot{\theta}
\end{equation}

Finally, the compact form is: \begin{equation}
\boxed{ \left( I + m l^2 \right) \ddot{\theta} + m g l \sin \theta 
= - m \ddot{x} \, (l \cos \theta) }
\end{equation}
\end{block}
\end{block}

\begin{block}{Linearization of the System}
\protect\phantomsection\label{linearization-of-the-system}
Since the analysis and control design techniques are only be done here
for linear systems, this set of equations needs to be linearized. About
the vertically upward equilibrium position, \(\theta = \pi\), and
assuming that the system stays within a small neighborhood of this
equilibrium, which is valid since deviation less than \(20^\circ\) from
the vertically upward position is desirable.

Let \(\phi\) represent the deviation of the pendulum's position from
equilibrium, that is, \(\theta = \pi + \phi\). Again presuming a small
deviation (\(\phi\)) from equilibrium, we can use the following small
angle approximations of the nonlinear functions in our system equations:

\begin{equation}
\cos \theta = \cos(\pi + \phi) \approx -1
\end{equation}

\begin{equation}
\sin \theta = \sin(\pi + \phi) \approx -\phi
\end{equation}

\begin{equation}
\dot{\theta}^2 = \dot{\phi}^2 \approx 0
\end{equation}

After substituting the above approximations into our nonlinear governing
equations, we arrive at the two linearized equations of motion. Note
\(u\) has been substituted for the input \(F\).

\begin{equation}
(I+ml^2)\ddot{\phi} - mgl\phi = ml\ddot{x}
\end{equation}

\begin{equation}
(M+m)\ddot{x} + b\dot{x} - ml\ddot{\phi} = u
\end{equation}
\end{block}

\begin{block}{Transfer Function}
\protect\phantomsection\label{transfer-function}
To obtain the transfer functions, we must first take the Laplace
transform of the system equations assuming zero initial conditions.

\begin{equation}
(I+ml^2)\Phi(s)s^2 - mgl\Phi(s) = mlX(s)s^2
\end{equation}

\begin{equation}
(M+m)X(s)s^2 + bX(s)s - ml\Phi(s)s^2 = U(s)
\end{equation}

Recall that a transfer function represents the relationship between a
single input and a single output at a time. To find our first transfer
function for the output \(\Phi(s)\) and an input of \(U(s)\) we need to
eliminate \(X(s)\) from the above equations. Solve the first equation
for \(X(s)\):

\begin{equation}
X(s) = \left[\frac{I+ml^2}{ml} - \frac{g}{s^2}\right]\Phi(s)
\end{equation}

Then substitute the above into the second equation:

\begin{equation}
(M+m)\left[\frac{I+ml^2}{ml} - \frac{g}{s^2}\right]\Phi(s)s^2 
+ b\left[\frac{I+ml^2}{ml} - \frac{g}{s^2}\right]\Phi(s)s 
- ml\Phi(s)s^2 = U(s)
\end{equation}

Rearranging, the transfer function is then the following:

\begin{equation}
\frac{\Phi(s)}{U(s)} = \frac{\tfrac{ml}{q}s^2}{s^4 + \tfrac{b(I+ml^2)}{q}s^3 - \tfrac{(M+m)mgl}{q}s^2 - \tfrac{bmgl}{q}s}
\end{equation}

where,

\begin{equation}
q = \big[(M+m)(I+ml^2) - (ml)^2\big]
\end{equation}

From the transfer function above it can be seen that there is both a
pole and a zero at the origin. These can be canceled and the transfer
function becomes the following:

\begin{equation}
P_{pend}(s) = \frac{\Phi(s)}{U(s)} = 
\frac{\tfrac{ml}{q}s}{s^3 + \tfrac{b(I+ml^2)}{q}s^2 - \tfrac{(M+m)mgl}{q}s - \tfrac{bmgl}{q}}
\qquad \left[\frac{\text{rad}}{\text{N}}\right]
\end{equation}

Second, the transfer function with the cart position \(X(s)\) as the
output can be derived in a similar manner to arrive at the following:

\begin{equation}
P_{cart}(s) = \frac{X(s)}{U(s)} = 
\frac{ \tfrac{ (I+ml^2)s^2 - gml } {q}}
{s^4 + \tfrac{b(I+ml^2)}{q}s^3 - \tfrac{(M+m)mgl}{q}s^2 - \tfrac{bmgl}{q}s}
\qquad \left[\frac{\text{m}}{\text{N}}\right]
\end{equation}
\end{block}

\begin{block}{State-Space Representation}
\protect\phantomsection\label{state-space-representation}
The linearized equations of motion from above can also be represented in
state-space form .

\begin{equation}
\begin{bmatrix}
  \dot{x}\\
  \ddot{x}\\
  \dot{\phi}\\
  \ddot{\phi}
\end{bmatrix}
=
\begin{bmatrix}
  0 & 1 & 0 & 0 \\
  0 & \tfrac{-(I+ml^2)b}{I(M+m)+Mml^2} & \tfrac{m^2gl^2}{I(M+m)+Mml^2} & 0 \\
  0 & 0 & 0 & 1 \\
  0 & \tfrac{-mlb}{I(M+m)+Mml^2} & \tfrac{mgl(M+m)}{I(M+m)+Mml^2} & 0
\end{bmatrix}
\begin{bmatrix}
  x \\
  \dot{x} \\
  \phi \\
  \dot{\phi}
\end{bmatrix}
+
\begin{bmatrix}
  0 \\
  \tfrac{I+ml^2}{I(M+m)+Mml^2} \\
  0 \\
  \tfrac{ml}{I(M+m)+Mml^2}
\end{bmatrix} u
\end{equation}

\begin{equation}
\mathbf{y} =
\begin{bmatrix}
  1 & 0 & 0 & 0 \\
  0 & 0 & 1 & 0
\end{bmatrix}
\begin{bmatrix}
  x \\
  \dot{x} \\
  \phi \\
  \dot{\phi}
\end{bmatrix}
+
\begin{bmatrix}
  0 \\
  0
\end{bmatrix} u
\end{equation}

The \(C\) matrix has 2 rows because both the cart's position and the
pendulum's position are part of the output. Specifically, the cart's
position is the first element of the output \(\mathbf{y}\) and the
pendulum's deviation from its equilibrium position is the second element
of \(\mathbf{y}\).
\end{block}
\end{frame}

\begin{frame}{MATLAB representation}
\protect\phantomsection\label{matlab-representation}
\begin{block}{Tranfer Function}
\protect\phantomsection\label{tranfer-function}
We can represent the transfer functions derived above for the inverted
pendulum system within MATLAB employing the following commands.

\[\text{From input } u \text{ to output } x:\] \[P_{cart}(s) = 
\frac{4.182 \times 10^{-6} s^{2} - 1.025 \times 10^{-4}}
     {2.3 \times 10^{-6} s^{4} + 4.182 \times 10^{-7} s^{3} - 7.172 \times 10^{-5} s^{2} - 1.025 \times 10^{-5} s}\]

\[\text{From input } u \text{ to output } \phi:\] \[P_{pend}(s) =
\frac{1.045 \times 10^{-5} s}
     {2.3 \times 10^{-6} s^{3} + 4.182 \times 10^{-7} s^{2} - 7.172 \times 10^{-5} s - 1.025 \times 10^{-5}}\]

\[\text{Continuous-time transfer function.}\]
\end{block}

\begin{block}{State -Space}
\protect\phantomsection\label{state--space}
We can also represent the system using the state-space equations. The
following additional MATLAB commands create a state-space model of the
inverted pendulum and produce the output shown below when run in the
MATLAB command window.

\begin{equation}
A = 
\begin{bmatrix}
0      & 1       & 0     & 0 \\
0      & -0.1818 & 2.673 & 0 \\
0      & 0       & 0     & 1 \\
0      & -0.4545 & 31.18 & 0
\end{bmatrix}, \quad
B =
\begin{bmatrix}
0 \\ 1.818 \\ 0 \\ 4.545
\end{bmatrix},
\end{equation}

\begin{equation}
C =
\begin{bmatrix}
1 & 0 & 0 & 0 \\
0 & 0 & 1 & 0
\end{bmatrix}, \quad
D =
\begin{bmatrix}
0 \\ 0
\end{bmatrix}
\end{equation}
\end{block}
\end{frame}

\begin{frame}[fragile]{Data Generation Analysis using Gymnasium}
\protect\phantomsection\label{data-generation-analysis-using-gymnasium}
This work uses the \texttt{InvertedPendulum-v5} environment from the
\texttt{Gymnasium} library, which is a MuJoCo-based environment. The
task is to balance an inverted pendulum on a cart by applying horizontal
forces.

\begin{block}{Import}
\protect\phantomsection\label{import}
\begin{verbatim}
import gymnasium as gym
env = gym.make("InvertedPendulum-v5")
\end{verbatim}
\end{block}

\begin{block}{Action Space}
\protect\phantomsection\label{action-space}
The action space is continuous and represented as:
\[\text{Box}([-3.0], [3.0], (1,), \texttt{float32})\]

The agent takes a one-dimensional continuous action (force applied to
the cart).

\begin{longtable}[]{@{}ccccccc@{}}
\caption{Action space specification of
InvertedPendulum-v5.}\tabularnewline
\toprule\noalign{}
Num & Action & Control Min & Control Max & Name & Joint & Type (Unit) \\
\midrule\noalign{}
\endfirsthead
\toprule\noalign{}
Num & Action & Control Min & Control Max & Name & Joint & Type (Unit) \\
\midrule\noalign{}
\endhead
0 & Force on cart & -3 & 3 & slider & slide & Force (N) \\
\bottomrule\noalign{}
\end{longtable}
\end{block}

\begin{block}{Observation Space}
\protect\phantomsection\label{observation-space}
The observation space is a 4-dimensional vector:
\[\text{Box}(-\infty, \infty, (4,), \texttt{float64})\]

It consists of cart position/velocity and pole angle/angular velocity.

\begin{longtable}[]{@{}ccccccc@{}}
\caption{Observation space specification of
InvertedPendulum-v5.}\tabularnewline
\toprule\noalign{}
Num & Observation & Min & Max & Name & Joint & Type (Unit) \\
\midrule\noalign{}
\endfirsthead
\toprule\noalign{}
Num & Observation & Min & Max & Name & Joint & Type (Unit) \\
\midrule\noalign{}
\endhead
0 & Cart position & \(-\infty\) & \(\infty\) & slider & slide & position
(m) \\
1 & Pole angle & \(-\infty\) & \(\infty\) & hinge & hinge & angle
(rad) \\
2 & Cart velocity & \(-\infty\) & \(\infty\) & slider & slide & velocity
(m/s) \\
3 & Pole angular velocity & \(-\infty\) & \(\infty\) & hinge & hinge &
angular velocity (rad/s) \\
\bottomrule\noalign{}
\end{longtable}
\end{block}

\begin{block}{Rewards}
\protect\phantomsection\label{rewards}
A reward of \(+1\) is given for each timestep that the pole remains
upright (within an angle limit).
\end{block}

\begin{block}{Starting State}
\protect\phantomsection\label{starting-state}
The initial state (position and velocity) is sampled from a multivariate
uniform continuous distribution, perturbed by random noise.
\end{block}

\begin{block}{Episode End}
\protect\phantomsection\label{episode-end}
\begin{itemize}
\item
  \textbf{Termination:} Episode terminates if:

  \begin{itemize}
  \item
    Any state value is non-finite.
  \item
    The absolute pole angle \(|\theta| > 0.2\) radians.
  \end{itemize}
\item
  \textbf{Truncation:} Default maximum of 1000 timesteps.
\end{itemize}
\end{block}

\begin{block}{Data Generation Code}
\protect\phantomsection\label{data-generation-code}
The Python code generates data from the \texttt{InvertedPendulum-v5}
environment using Gymnasium. It collects state transitions, actions, and
rewards for further analysis.

Here we are considering the agent to take completely random actions:

\textbf{Main goal:} The script is for \textbf{generating datasets} the
Inverted Pendulum system. It runs the simulation for 50 episodes and
saves the detailed step-by-step results of each run into separate
\texttt{.csv} files. These datasets serve as a \textbf{baseline
performance} for a random agent, which can later be used to evaluate the
performance of a ``smart'' agent.

\begin{block}{How the Code Works}
\protect\phantomsection\label{how-the-code-works}
The script simulates an inverted pendulum by interacting with a
\texttt{gym} environment and collecting data. Its workflow is as
follows:

\begin{enumerate}
\item
  \textbf{Environment Setup:} The environment is created using
  \texttt{env\ =\ gym.make(...)} and initialized with
  \texttt{env.reset()}.
\item
  \textbf{Episode Loop:} For each episode, the pendulum is simulated
  until it either falls (\texttt{terminated}) or the time limit is
  reached (\texttt{truncated}).
\item
  \textbf{Step Loop:} Within each episode, random actions are generated
  (\texttt{env.action\_space.sample()}) and applied
  (\texttt{env.step(action)}). The resulting next state, reward, and
  episode status are returned.
\item
  \textbf{Data Collection:} At each step, the current state, chosen
  action, reward, next state, and termination flag are recorded in
  \texttt{collected\_data}. The current observation is updated
  accordingly.
\item
  \textbf{Cleanup:} After all episodes, the environment is closed using
  \texttt{env.close()}.
\end{enumerate}
\end{block}
\end{block}

\begin{block}{Analysis of Simulation Plots}
\protect\phantomsection\label{analysis-of-simulation-plots}
The following plots were generated for (\(N = 50\)) numbers of episodes
. Each plot provides a different perspective on the behavior of the
inverted pendulum under random actions.

\begin{block}{Histogram of Episode Durations}
\protect\phantomsection\label{histogram-of-episode-durations}
\begin{figure}[H]
\centering
\includegraphics[width=0.7\linewidth,height=\textheight,keepaspectratio]{plots/episode_durations_histogram_50.png}
\caption{Histogram of episode durations for
\(N=50\).}\label{fig:histogram}
\end{figure}

This plot serves as the main performance indicator, showing how
successful the random agent was at balancing the pole.

\begin{itemize}
\item
  \textbf{X-axis}: Episode duration (number of steps).
\item
  \textbf{Y-axis}: Frequency of episodes with that duration.
\item
  \textbf{Interpretation}: For a random agent, the distribution is
  typically skewed to the left, meaning most episodes are very short. A
  few ``lucky'' runs may last longer, forming a long right tail. With
  larger \(N\), the histogram gives a more reliable picture of
  performance.
\item
  \textbf{Insight:} This confirms that the agent's performance is
  consistently poor, failing almost immediately in most attempts.
\end{itemize}
\end{block}

\begin{block}{Phase Plot: Pole Angle vs. Angular Velocity}
\protect\phantomsection\label{phase-plot-pole-angle-vs.-angular-velocity}
\begin{figure}[H]
\centering
\includegraphics[width=0.7\linewidth,height=\textheight,keepaspectratio]{plots/phase_plot_50.png}
\caption{Phase plot: Pole angle vs. angular velocity.}\label{fig:phase}
\end{figure}

This phase-space plot captures the underlying physical dynamics of the
pendulum.

\begin{itemize}
\item
  \textbf{X-axis}: Pole angle (radians), where \(0\) is upright.
\item
  \textbf{Y-axis}: Angular velocity of the pole (rad/s).
\item
  \textbf{Interpretation}: Each point represents the system state at a
  particular timestep. The goal state is \((0,0)\), meaning upright and
  motionless. For a random policy, the points typically spread outward
  from the center, reflecting that once the pole starts to fall, its
  angular velocity grows quickly, leading to failure.
\item
  \textbf{Insight:} This illustrates the inherent instability of the
  inverted pendulum. Random actions cannot restore the system state back
  toward the stable center.
\end{itemize}
\end{block}

\begin{block}{Time Series of the Longest Episode}
\protect\phantomsection\label{time-series-of-the-longest-episode}
\begin{figure}[H]
\centering
\includegraphics[width=0.7\linewidth,height=\textheight,keepaspectratio]{plots/longest_episode_timeseries_50.png}
\caption{Time series of the longest episode.}\label{fig:time_series}
\end{figure}

This plot highlights the trajectory of the \emph{most successful}
episode (the one who succeeded with the maximum number of timesteps
without failing).

\begin{itemize}
\item
  \textbf{X-axis}: Timestep within the episode.
\item
  \textbf{Y-axis}: Values of system variables (cart position, pole
  angle, etc.).
\item
  \textbf{Interpretation}: The plot shows how state variables evolve
  over time, alongside the (random) control actions. It gives insight
  into how actions affect the system dynamics, even if the control
  policy is not optimal.
\item
  \textbf{Insight:} The pole angle stays close to zero for a short
  period before drifting too far, causing the episode to end. The action
  signal is erratic, even in this ``successful'' run. This shows that
  the agent has no control strategy---only chance kept the system
  balanced briefly.
\end{itemize}
\end{block}

\begin{block}{Reward per Episode}
\protect\phantomsection\label{reward-per-episode}
\begin{figure}[H]
\centering
\includegraphics[width=0.7\linewidth,height=\textheight,keepaspectratio]{plots/rewards_per_episode_50.png}
\caption{Total reward earned in each episode.}\label{fig:reward}
\end{figure}

This bar chart summarizes the total reward earned in each episode.

\begin{itemize}
\item
  \textbf{X-axis}: Episode index.
\item
  \textbf{Y-axis}: Total reward (equivalent to episode duration).
\item
  \textbf{Interpretation}: With random actions, episode rewards
  fluctuate without any trend. Most bars are short (early failures),
  with a few taller bars (lucky long runs). A red dashed line indicates
  the average reward across episodes, which remains low for a random
  agent.
\item
  \textbf{Insight:} The reward (i.e., duration) fluctuates randomly
  across the \(50\) episodes, with no upward trend. The average reward
  (red dashed line) is very low. This reinforces that the agent is not
  learning or improving. Its success is erratic and depends entirely on
  luck.
\end{itemize}
\end{block}
\end{block}
\end{frame}

\begin{frame}[fragile]{Conclusion and Improvements}
\protect\phantomsection\label{conclusion-and-improvements}
\begin{block}{Conclusion about Random Agent}
\protect\phantomsection\label{conclusion-about-random-agent}
The four plots paint a clear and consistent picture: a random action
strategy is completely inadequate for solving the Inverted Pendulum
problem. The agent exhibits no learning or strategic behavior, and its
performance relies purely on chance. These results establish a
quantitative baseline for failure, against which more advanced control
and reinforcement learning approaches can be compared.
\end{block}

\begin{block}{Improvements}
\protect\phantomsection\label{improvements}
Here we outline one of the potential improvements that can be employed
to achieve meaningful control of the Inverted Pendulum.

The current random agent provides a baseline, but it cannot learn or
adapt. The following improvements describe how to create a more
intelligent agent that learns to balance the Inverted Pendulum
effectively.

\begin{block}{Implement a Reinforcement Learning Algorithm}
\protect\phantomsection\label{implement-a-reinforcement-learning-algorithm}
The fundamental change is to replace the random action selection with a
learning algorithm. Instead of guessing, the agent will have a
\emph{policy} that it improves over time. For the Inverted Pendulum
problem, the recommended approach is to use a Reinforcement Learning
(RL) algorithm, which is \textbf{Proximal Policy Optimization (PPO)} in
this case.

PPO is a powerful, stable, and widely used RL algorithm. It works
exceptionally well for control problems like the Inverted Pendulum and
is a great starting point to achieve strong results.
\end{block}

\begin{block}{Use a Reinforcement Learning Library}
\protect\phantomsection\label{use-a-reinforcement-learning-library}
Implementing PPO from scratch is complex. Hence we are using a pre-built
library that handles the low-level details, \textbf{Library: Stable
Baselines3} which integrates seamlessly with Gymnasium environments and
simplifies training a powerful agent.
\end{block}

\begin{block}{Train the Agent}
\protect\phantomsection\label{train-the-agent}
\begin{enumerate}
\item
  \textbf{Training Phase} In this phase, the agent interacts with the
  environment to learn a policy. The Stable Baselines3 library handles
  the underlying RL loop:

  \begin{itemize}
  \item
    The agent interacts with the environment to collect experience.
  \item
    It uses that experience to update its policy.
  \item
    This process repeats for a specified number of steps (e.g., 10k,
    20k), progressively improving the agent's ability to balance the
    pole.
  \end{itemize}
\item
  \textbf{Evaluation Phase}

  After training, the saved trained model is used to control the
  environment:

  \begin{itemize}
  \item
    Instead of choosing random actions, the trained model predicts the
    optimal action at each step.
  \item
    The agent can now balance the pole for much longer durations, often
    reaching the full episode limit (e.g., 1,000 steps), demonstrating
    learned control.
  \end{itemize}
\end{enumerate}

By switching from random actions to a PPO-based policy, the agent
transitions from purely chance-driven behavior to intelligent,
goal-directed control.
\end{block}
\end{block}

\begin{block}{Analysis of Simulation Plots}
\protect\phantomsection\label{analysis-of-simulation-plots-1}
The following plots were generated for (\(N=50\)) numbers of episodes.
Each plot provides a different perspective on the behavior of the
inverted pendulum under random actions.

\begin{block}{Histogram of Episode Durations}
\protect\phantomsection\label{histogram-of-episode-durations-1}
\begin{figure}[H]
\centering
\includegraphics[width=0.7\linewidth,height=\textheight,keepaspectratio]{plots/trained_episode_durations_histogram_50.png}
\caption{Histogram of episode durations for
\(N=50\).}\label{fig:histogram}
\end{figure}

This plot serves as the main performance indicator, showing how
successful the random agent was at balancing the pole.

\begin{itemize}
\item
  \textbf{Observation:} The histogram of episode durations exhibits a
  bimodal distribution. A good number of the episodes successfully reach
  the 1,000-step truncation limit, as indicated by the prominent peak on
  the far right. However, other cluster of episodes terminate
  prematurely at various shorter durations.
\item
  \textbf{Conclusion:} This distribution signifies a policy that is
  successful under most tested conditions but is not yet robust enough
  to generalize to all possible states. The agent has clearly learned
  the primary stabilization task, but its strategy fails when
  encountering certain less-frequently experienced or more challenging
  state-space regions.
\end{itemize}
\end{block}

\begin{block}{Phase Plot: Pole Angle vs. Angular Velocity}
\protect\phantomsection\label{phase-plot-pole-angle-vs.-angular-velocity-1}
\begin{figure}[H]
\centering
\includegraphics[width=0.7\linewidth,height=\textheight,keepaspectratio]{plots/trained_phase_plot_50.png}
\caption{Phase plot: Pole angle vs. angular velocity.}\label{fig:phase}
\end{figure}

This phase-space plot captures the underlying physical dynamics of the
pendulum.

\begin{itemize}
\item
  \textbf{Observation:} The phase points are tightly clustered around
  the origin \((0,0)\), very similar to the perfectly trained case.
\item
  \textbf{Conclusion:} This indicates that the agent maintains the
  pendulum in a balanced and stable state for most of the time. The rare
  failures are too brief to affect this plot significantly.
\end{itemize}
\end{block}

\begin{block}{Time Series of the Longest Episode}
\protect\phantomsection\label{time-series-of-the-longest-episode-1}
\begin{figure}[H]
\centering
\includegraphics[width=0.7\linewidth,height=\textheight,keepaspectratio]{plots/trained_longest_episode_timeseries_50.png}
\caption{Time series of the longest episode.}\label{fig:time_series}
\end{figure}

This plot highlights the trajectory of the \emph{most successful}
episode.

\begin{itemize}
\item
  \textbf{Observation:} The 1000-step run demonstrates that the pole's
  angle is controlled effectively, with smooth and deliberate actions
  rather than random movements.
\item
  \textbf{Conclusion:} This confirms that the agent has indeed learned
  the \emph{correct control strategy}. When it succeeds, it does so with
  skill, not by chance.
\end{itemize}
\end{block}

\begin{block}{Reward per Episode}
\protect\phantomsection\label{reward-per-episode-1}
\begin{figure}[H]
\centering
\includegraphics[width=0.7\linewidth,height=\textheight,keepaspectratio]{plots/trained_rewards_per_episode_50.png}
\caption{Total reward earned in each episode.}\label{fig:reward}
\end{figure}

This bar chart summarizes the total reward earned in each episode.

\begin{itemize}
\item
  \textbf{Observation:} The bar chart displays the total reward for each
  of the 50 evaluation episodes. The vast majority of episodes achieve
  the maximum possible reward of 1000, as shown by the consistently tall
  bars. However, a small number of episodes terminate prematurely,
  resulting in lower rewards.
\item
  \textbf{Insight:} The agent has successfully learned a highly
  effective control policy, transitioning from the chaotic,
  low-performance behavior of a random agent to a skilled, strategic
  one.
\end{itemize}
\end{block}
\end{block}

\begin{block}{Final Conclusion}
\protect\phantomsection\label{final-conclusion}
\begin{itemize}
\item
  The agent trained for \(10{,}000\) timesteps demonstrates a
  fundamentally different and vastly superior mode of operation compared
  to the baseline random policy.
\item
  The random policy is an open-loop strategy, where actions are
  independent of the system state. This leads to:

  \begin{itemize}
  \item
    rapid system divergence,
  \item
    uniformly poor performance,
  \item
    complete inability to stabilize the pendulum,
  \item
    universally short episode durations.
  \end{itemize}
\item
  In contrast, the trained policy is a closed-loop control strategy that
  maps states to actions effectively. Its advantages include:

  \begin{itemize}
  \item
    \textbf{Achieving Stability:} actively drives the system state
    towards the stable equilibrium point, unlike the random policy which
    accelerates divergence.
  \item
    \textbf{Sustaining Performance:} balances the pendulum up to the
    maximum time limit in most trials, a feat statistically impossible
    for a random policy.
  \item
    \textbf{Exhibiting Skillful Control:} actions are deliberate and
    corrective, opposing the system's unstable dynamics, unlike the
    uncorrelated and ineffective actions of the random policy.
  \end{itemize}
\end{itemize}

Furthermore, increasing the training horizon, i.e., the
\texttt{total\_timesteps}, enhances the

agent's performance. With more training iterations, the agent converges
to more accurate

policies, yielding greater stability, robustness, and improved control
strategies compared

to those obtained with fewer training steps.
\end{block}
\end{frame}
